\chapter{Esercizi: Array in Java}
\label{chap:esercizi-array-java}

\begin{exercise}
\label{ex:fi-esercizi-array-funzione}
Che problema risolve un array quando si deve memorizzare una sequenza
di molti valori dello stesso tipo?
\end{exercise}

Permette di accedere ai valori mediante un riferimento e un indice,
senza dichiarare una variabile distinta per ciascun elemento
(\cref{def:fi-java-array}).

\begin{exercise}
\label{ex:fi-esercizi-array-default}
Quali valori iniziali assumono automaticamente gli elementi di un
array di \texttt{int}, di \texttt{boolean} e di riferimenti a oggetti?
\end{exercise}

Rispettivamente \texttt{0}, \texttt{false} e \texttt{null}
(\cref{sec:fi-java-array-inizializzazione}).

\begin{exercise}
\label{ex:fi-esercizi-array-indici}
Che cosa succede se si accede a un array con un indice negativo
oppure con un indice uguale alla lunghezza dell'array?
\end{exercise}

L'accesso provoca \texttt{ArrayIndexOutOfBoundsException}:
gli indici validi vanno da zero a \texttt{length - 1}
(\cref{sec:fi-java-array-costruzione}).

\begin{exercise}
\label{ex:fi-esercizi-array-length}
Perché \texttt{values.length} non è un metodo, e perché non è
possibile assegnare un nuovo valore a \texttt{values.length}?
\end{exercise}

\texttt{length} è un campo pubblico e \texttt{final} dell'oggetto
array, non un metodo. La lunghezza è fissata alla costruzione
e l'assegnamento al campo è vietato dal compilatore
(\cref{sec:fi-java-array-inizializzazione}).

\begin{exercise}
\label{ex:fi-esercizi-array-validazione}
Come si può controllare che un indice inserito dall'utente sia
valido prima di usare \texttt{values[index]}?
\end{exercise}

Si verifica \texttt{index >= 0 \&\& index < values.length},
dopo aver accertato che \texttt{values} non sia \texttt{null}
(\cref{sec:fi-java-array-costruzione}).

\begin{exercise}
\label{ex:fi-esercizi-array-inizializzazione}
Quali sono i tre modi mostrati nelle slide per inizializzare un array
con i valori \texttt{2}, \texttt{3}, \texttt{5}, e quale forma è
preferibile quando i valori sono già noti?
\end{exercise}

Si possono costruire tre celle con \texttt{new int[3]} e assegnare
singolarmente gli elementi; usare \texttt{int[] a = \{2, 3, 5\};};
oppure \texttt{int[] a = new int[] \{2, 3, 5\};}.
Nella dichiarazione, la seconda forma è la più concisa
(\cref{sec:fi-java-array-inizializzazione}).

\begin{exercise}
\label{ex:fi-esercizi-array-null}
Perché un metodo che riceve un array dovrebbe spesso controllare
se il parametro vale \texttt{null}?
\end{exercise}

Un riferimento nullo non individua un array.
L'accesso a \texttt{v.length} in questo caso provoca
\texttt{NullPointerException}. Una precondizione esplicita
consente di segnalare l'argomento non valido (\cref{sec:fi-java-array-parametri}).

\begin{exercise}
\label{ex:fi-esercizi-array-ritorno}
Come può un metodo restituire un array come valore di ritorno,
per esempio nel caso di un metodo \texttt{resize}?
\end{exercise}

Dichiara un tipo di ritorno array, come \texttt{int[]}, e
restituisce il riferimento all'array costruito. Il chiamante può
assegnare quel riferimento alla propria variabile
(\cref{sec:fi-java-array-resize}).

\begin{exercise}
\label{ex:fi-esercizi-array-alias}
Perché l'assegnazione \texttt{double[] y = x;} non copia il
contenuto dell'array, ma solo il riferimento allo stesso oggetto?
\end{exercise}

Le variabili di tipo array contengono riferimenti. L'assegnamento
copia il valore del riferimento, quindi entrambe le variabili
individuano lo stesso array
(\cref{sec:fi-java-array-parametri}).

\begin{exercise}
\label{ex:fi-esercizi-array-copia}
Quali passi sono necessari per ottenere una vera copia
indipendente di un array?
\end{exercise}

Si costruisce un secondo array e si copiano gli elementi, usando
un ciclo, \texttt{System.arraycopy()} o \texttt{clone()}.
Se gli elementi sono riferimenti, gli oggetti indicati
rimangono condivisi (\cref{sec:fi-java-array-parametri}).

\begin{exercise}
\label{ex:fi-esercizi-array-arraycopy}
A che cosa serve \texttt{System.arraycopy} e quale significato
hanno i parametri \texttt{from}, \texttt{fromStart}, \texttt{to},
\texttt{toStart} e \texttt{count}?
\end{exercise}

Copia \texttt{count} elementi dall'array \texttt{from}, a partire
da \texttt{fromStart}, nell'array \texttt{to}, a partire da
\texttt{toStart}. Gli intervalli devono essere validi e i tipi
compatibili (\cref{sec:fi-java-array-parametri}).

\begin{exercise}
\label{ex:fi-esercizi-array-clone}
Perché il metodo \texttt{clone} può essere usato per copiare
un array, e perché è necessario effettuare un cast?
\end{exercise}

\texttt{clone()} crea un nuovo array con gli stessi elementi.
Per gli array Java il cast non è necessario: il tipo restituito
è quello dell'array, ad esempio \texttt{double[]}, anche se
\texttt{Object.clone()} in generale restituisce \texttt{Object}.
La copia rimane superficiale
(\cref{sec:fi-java-array-parametri}).

\begin{exercise}
\label{ex:fi-esercizi-array-foreach}
Che cos'è il ciclo \texttt{for-each} e perché riduce il rischio
di errori legati agli indici?
\end{exercise}

Visita direttamente gli elementi di un array o di una collezione,
senza controllare esplicitamente gli indici; evita quindi errori
nelle condizioni di avanzamento
(\cref{sec:fi-java-array-for-each}).

\begin{exercise}
\label{ex:fi-esercizi-array-foreach-limiti}
In quale situazione un ciclo \texttt{for-each} è meno adatto
di un ciclo \texttt{for} tradizionale?
\end{exercise}

Quando occorre conoscere l'indice, visitare solo alcune posizioni
o sostituire le celle dell'array durante la scansione
(\cref{sec:fi-java-array-for-each}).

\begin{exercise}
\label{ex:fi-esercizi-array-ridimensionamento}
Come funziona il metodo \texttt{resize} usato per creare un
array più grande?
\end{exercise}

Costruisce un nuovo array della lunghezza richiesta, copia
gli elementi conservabili e restituisce il nuovo riferimento.
L'array originario non cambia dimensione
(\cref{sec:fi-java-array-resize}).

\begin{exercise}
\label{ex:fi-esercizi-array-stampa}
Perché stampare direttamente un array con
\texttt{System.out.println(a)} non mostra il contenuto dell'array?
\end{exercise}

L'array non ridefinisce una rappresentazione testuale dei propri
elementi. La stampa usa la rappresentazione dell'oggetto;
per mostrare i valori occorre elaborarli oppure utilizzare
un metodo che li converta in testo
(\cref{sec:fi-java-array-stampa}).

\begin{exercise}
\label{ex:fi-esercizi-array-inserimento}
Perché nell'inserimento di un elemento in mezzo a un array
gli spostamenti devono essere eseguiti dal basso verso l'alto?
\end{exercise}

Si parte dall'ultimo elemento e si procede verso il punto di
inserimento. In questo modo ciascun valore viene spostato
nell'indice successivo prima che la sua posizione originale
sia sovrascritta (\cref{sec:fi-java-array-insert}).

\begin{exercise}
\label{ex:fi-esercizi-array-paralleli}
Perché è preferibile sostituire gli array paralleli con un
array di oggetti?
\end{exercise}

Una classe raccoglie le proprietà dello stesso elemento logico.
L'aggiunta o la modifica di un dato non impone la gestione
coordinata di più array o l'estensione delle firme di numerosi
metodi (\cref{sec:fi-java-array-paralleli}).
